\begin{afterword}
\summaryhead

The post poses two choices: \$24,000 for certain against a 33/34 chance of \$27,000, and a 34\% chance of \$24,000 against a 33\% chance of \$27,000. Most people take the sure thing in the first and the larger prize in the second, and a majority of subjects make both choices. Since each second gamble is the first played with probability 0.34, no utility function makes both choices maximize expected utility; the pattern breaks the independence axiom, whatever the utility of money.

Allais, who published the problem in 1953, took it to show a flaw in the theory. The author replies that intuitions reveal how the mind works, not directly which choices are wise, and that ``There is always a price for leaving the Bayesian Way.'' The price here is a money pump: in a two-stage game an agent with both preferences pays a penny to switch to the gamble before a die roll and another penny to switch back after it.

\responsehead

The mathematical core of the post is correct. The two inequalities cannot both hold for any utility function, the arithmetic of the two-stage game is right, and the claim that a majority of subjects make both choices matches the paper the post lists: 61 per cent in Kahneman and Tversky's version. The post's own numbers change the problem from Allais's best-known form to a related one, but both break the same axiom.

What the post adds to this is a verdict: the pattern is an error in the subjects, not a flaw in the theory. That verdict is one side of a dispute that is still open. On the post's side is Savage, who argued in 1954 that the Allais preferences are irrational and reported that he had held them himself before changing his mind. On the other side is Allais, who in 1979 published a long rejoinder extending his original view, and recent theorists such as Buchak, who argue that a rational agent may weigh risk itself. The post treats the opposing view as a feeling of reassurance, ``naive philosophical realism.'' Its published defenders also give theories that keep a consistent ranking of gambles and drop only independence. The post's ``consistent'' means consistent with expected utility, and that is the point at issue.

The money pump is the post's argument for its verdict, and ``There is always a price'' states it as settled. The pump takes two pennies from an agent who does not look ahead. The decision-theory literature asks what happens to one who does. On Hammond's argument, which supports the post, an agent who plans ahead would foresee switching back and end with the gamble it likes less; defenders of independence-breaking theories reply that such an agent can adopt a plan and keep to it. Either way the agent loses at most one penny. The evidence adds a complication: when Kahneman and Tversky described a game of this shape in two stages, most subjects chose the sure option before the first stage as well, so the reversal the pump needs did not appear, and a dependence on description appeared instead.

In short: a correct proof that the common Allais choices cannot both maximize expected utility, followed by the verdict that they are a mistake, which is Savage's side of a dispute with Allais that the money pump, as described, does not settle.
\end{afterword}
